\documentclass[12pt, a4paper]{article}

\usepackage{geometry}
\geometry{
	a4paper,
	left=30mm,
	right=10mm,
	top=20mm,
	bottom=20mm
}

\usepackage{fontspec}
\usepackage[ukrainian, english, provide=*]{babel}

\babelfont{rm}{Times New Roman}
\babelfont{sf}{Arial}
\babelfont{tt}{Courier New}

\usepackage{setspace}
\onehalfspacing

\usepackage{indentfirst}
\setlength{\parindent}{1.25cm}
\setlength{\parskip}{0pt}

\usepackage{amsmath, amssymb}
\usepackage{booktabs}
\usepackage{tabularx}
\usepackage{array}
\usepackage{caption}
\captionsetup{labelsep=period}

\usepackage{hyperref}
\hypersetup{
	colorlinks=true,
	linkcolor=black,
	urlcolor=black,
	citecolor=black
}

\begin{document}
	\selectlanguage{ukrainian}
	
	% ==========================================
	% ТИТУЛЬНА СТОРІНКА
	% ==========================================
	\begin{titlepage}
		\centering
		{\fontsize{14pt}{16pt}\selectfont \textbf{МІНІСТЕРСТВО ОСВІТИ І НАУКИ УКРАЇНИ}}\\[0.3em]
		{\fontsize{14pt}{16pt}\selectfont \textbf{ОДЕСЬКИЙ НАЦІОНАЛЬНИЙ УНІВЕРСИТЕТ ІМЕНІ І. І. МЕЧНИКОВА}}\\[0.3em]
		{\fontsize{14pt}{16pt}\selectfont Факультет математики, фізики та інформаційних технологій}\\[0.3em]
		{\fontsize{14pt}{16pt}\selectfont Кафедра оптимального керування і економічної кібернетики}\\[5em]
		
		{\fontsize{20pt}{24pt}\selectfont \textbf{ЗВІТ}}\\[1.2em]
		{\fontsize{16pt}{20pt}\selectfont \textbf{з лабораторної роботи №4}}\\[0.6em]
		{\fontsize{14pt}{18pt}\selectfont з дисципліни «Системи комп'ютерного зору»}\\[1.2em]
		{\fontsize{14pt}{18pt}\selectfont \textbf{Тема:} «Зіставлення ознак, RANSAC, побудова панорами»}\\[6em]
		
		\begin{flushright}
			\fontsize{14pt}{18pt}\selectfont
			\textbf{Виконала:}\\[0.2em]
			студентка 3 курсу\\[0.2em]
			спеціальності «Прикладна математика»\\[0.2em]
			Овчарук Єлизавета Євгеніївна\\[2em]
			
			\textbf{Перевірив:}\\[0.2em]
			доцент, к. ф.-м. н.\\[0.2em]
			Мороз Володимир Володимирович
		\end{flushright}
		
		\vfill
		{\fontsize{14pt}{16pt}\selectfont Одеса --- 2026}
	\end{titlepage}
	
	\newpage
	
	% ==========================================
	% ВСТУП І МЕТА РОБОТИ
	% ==========================================
	\section*{Вступ і мета роботи}
	\addcontentsline{toc}{section}{Вступ і мета роботи}
	
	\textbf{Мета роботи:} Дослідити алгоритми зіставлення дескрипторів локальних ознак (SIFT), критерії відсіювання неоднозначних відповідностей (відношення відстаней Лоу, взаємна перевірка), вивчити проблеми класичного методу найменших квадратів (МНК) в умовах викидів, реалізувати власний робастний алгоритм RANSAC для оцінки гомографії, дослідити залежність якості від порога репроєкції та побудувати цілісну панораму серії кадрів із лінійним блендингом.
	
	\textbf{Завдання:}
	\begin{enumerate}
		\item Обчислити ключові точки й дескриптори SIFT для суміжних кадрів.
		\item Виконати зіставлення дескрипторів через $k$-NN ($k=2$) та перевірити роботу критерію Лоу для значень $T \in [0.5, 0.95]$.
		\item Застосувати взаємну симетричну перевірку (cross-check).
		\item Оцінити гомографію МНК без фільтрації та продемонструвати катастрофічний вплив викидів на похибку репроєкції.
		\item Реалізувати власний алгоритм RANSAC, розрахувати теоретичну кількість ітерацій $N$ та порівняти результат з реалізацією OpenCV.
		\item Дослідити поведінку метрик при зміні порога $\tau \in [1, 20]$ px і виявити плато стабільності.
		\item Побудувати панораму з 5 послідовних кадрів відносно опорного кадру та обчислити середню похибку репроєкції ланцюжка перетворень.
		\item Дослідити випадки відмови гомографічної моделі.
	\end{enumerate}
	
	% ==========================================
	% ТЕОРЕТИЧНІ ВІДОМОСТІ
	% ==========================================
	\section{Теоретичні відомості}
	
	\subsection{Дескриптори SIFT та критерій відношення Лоу}
	Алгоритм SIFT (Scale-Invariant Feature Transform) знаходить масштабно-інваріантні екстремуми в різницях гаусіанів (DoG) та будує для кожного околу 128-вимірний вектор градієнтних гістограм. Порівняння дескрипторів $\mathbf{d}_1, \mathbf{d}_2$ здійснюється за евклідовою нормою:
	\begin{equation}
		D(\mathbf{d}_1, \mathbf{d}_2) = \|\mathbf{d}_1 - \mathbf{d}_2\|_2 = \sqrt{\sum_{i=1}^{128} (d_{1,i} - d_{2,i})^2}
	\end{equation}
	
	Для відсіювання помилкових відповідностей Девід Лоу запропонував критерій відношення відстаней (Lowe's ratio test): шукаються перший $m_1$ та другий $m_2$ найближчі сусіди. Зіставлення вважається достовірним, якщо:
	\begin{equation}
		\frac{\|\mathbf{d} - \mathbf{d}_{m_1}\|}{\|\mathbf{d} - \mathbf{d}_{m_2}\|} < T, \quad T \in (0, 1)
	\end{equation}
	Якщо перший і другий сусіди розташовані на близькій відстані ($T \approx 1$), це сигналізує про регулярну структуру сцени або шум.
	
	\subsection{Проєктивна геометрія та матриця гомографії}
	Гомографія $H \in \mathbb{R}^{3 \times 3}$ описує проєктивне відображення точок площини між двома різними ракурсами камери або між довільними сценами при чистому обертанні камери навколо власного оптичного центру:
	\begin{equation}
		\mathbf{x}' \sim H \mathbf{x} \iff \begin{bmatrix} x' \\ y' \\ 1 \end{bmatrix} \sim \begin{bmatrix} h_{11} & h_{12} & h_{13} \\ h_{21} & h_{22} & h_{23} \\ h_{31} & h_{32} & h_{33} \end{bmatrix} \begin{bmatrix} x \\ y \\ 1 \end{bmatrix}
	\end{equation}
	Матриця визначена з точністю до масштабу і має 8 ступенів вільності, тому для її визначення потрібно щонайменше 4 неколінеарні пари відповідних точок.
	
	\subsection{Робастний алгоритм RANSAC}
	Метод найменших квадратів мінімізує суму квадратів відхилень, тому навіть поодинокі грубі викиди (outliers) різко зміщують оцінку. Алгоритм RANSAC (Random Sample Consensus) вирішує цю проблему:
	\begin{enumerate}
		\item Випадково обирається мінімальна підмножина точок ($s = 4$).
		\item Оцінюється пробна матриця $H$.
		\item Обчислюється похибка репроєкції для всіх $M$ точок:
		\begin{equation}
			d_i = \|\mathbf{x}'_i - \hat{\mathbf{x}}'_i\|_2 = \left\| \mathbf{x}'_i - \frac{1}{\mathbf{h}_3^T \mathbf{x}_i} \begin{bmatrix} \mathbf{h}_1^T \mathbf{x}_i \\ \mathbf{h}_2^T \mathbf{x}_i \end{bmatrix} \right\|_2
		\end{equation}
		\item Точка кваліфікується як інлаєр (узгоджена), якщо $d_i < \tau$.
		\item Модель із максимальною кількістю інлаєрів зберігається, і фінальна матриця $H$ переоцінюється за всіма виявленими інлаєрами.
	\end{enumerate}
	
	Теоретична кількість ітерацій $N$, необхідна для вибору хоча б однієї вибірки суто з правильних точок із заданою ймовірністю $p$ ($p=0.99$):
	\begin{equation}
		N = \frac{\ln(1 - p)}{\ln(1 - (1 - \varepsilon)^s)}
	\end{equation}
	де $\varepsilon = 1 - \frac{N_{\text{inliers}}}{N_{\text{total}}}$ --- спостережувана частка викидів.
	
	% ==========================================
	% ПРАКТИЧНЕ ВИКОНАННЯ ТА РЕЗУЛЬТАТИ
	% ==========================================
	\section{Практичне виконання та результати експериментів}
	
	Для дослідження використано набір тестових зображень набережної (\texttt{boat1.jpg} -- \texttt{boat5.jpg}).
	
	\subsection{Частина A. Зіставлення ознак}
	На перших двох кадрах детектовано:
	\begin{itemize}
		\item Кадр 1: 19\,096 ключових точок SIFT;
		\item Кадр 2: 12\,804 ключових точок SIFT.
	\end{itemize}
	
	Застосовано \texttt{cv2.BFMatcher(cv2.NORM\_L2)} із двостороннім $k$-NN пошуком ($k=2$). Досліджено вплив порога $T$ критерію Лоу на кількість кандидатів (табл.~\ref{tab:lowe}).
	
	\begin{table}[ht]
		\centering
		\caption{Залежність кількості кандидатів від порога $T$}
		\label{tab:lowe}
		\begin{tabular}{ccccccc}
			\toprule
			\textbf{Порог $T$} & 0.50 & 0.60 & 0.70 & 0.80 & 0.90 & 0.95 \\
			\midrule
			\textbf{Кількість пар} & 2117 & 2854 & 3490 & 4268 & 6593 & 10298 \\
			\bottomrule
		\end{tabular}
	\end{table}
	
	При значенні $T=0.75$ отримано 3839 кандидатів. Після застосування умови взаємної узгодженості (cross-check) залишилося 3485 надійних пар. Сирі відповідності до фільтрації містили значну кількість перехресних ліній (помилкових зв'язків між хвилями води та небом).
	
	\begin{figure}[ht]
		\centering
		\framebox{\parbox{0.92\textwidth}{\centering \textbf{Рис. 1.} Результат зіставлення сирих кандидатів (\texttt{results/step5\_raw\_matches.jpg}).\\
				\small Чітко видно численні перехресні та некоректні промені між водою, набережною та небом через наявність шумів та повторюваних текстур.}}
	\end{figure}
	
	\subsection{Частина B. Порівняння з методом найменших квадратів}
	Оцінка гомографії через стандартний МНК без відсіву (\texttt{method = 0}):
	\begin{equation*}
		H_{\text{LS}} = \begin{bmatrix}
			9.337 \cdot 10^{-1} & -7.874 \cdot 10^{-2} & -8.721 \cdot 10^{2} \\
			-1.917 \cdot 10^{-2} & 1.190 \cdot 10^{0} & -1.868 \cdot 10^{2} \\
			-2.155 \cdot 10^{-5} & 4.409 \cdot 10^{-5} & 1.000 \cdot 10^{0}
		\end{bmatrix}
	\end{equation*}
	
	Середня похибка репроєкції для МНК склала:
	\begin{equation*}
		\text{Error}_{\text{LS}} = 59.64 \text{ px}
	\end{equation*}
	Трансформоване зображення зазнало критичного геометричного зламу та перекосу горизонту через наявність $\sim 12\%$ викидів, які повністю змістили оптимальну площину.
	
	\begin{figure}[ht]
		\centering
		\framebox{\parbox{0.92\textwidth}{\centering \textbf{Рис. 2.} Деформація першого кадру за матрицею гомографії МНК (\texttt{results/step7\_warped\_least\_squares.jpg}).\\
				\small Критичний завал площини та спотворення геометрії через чутливість квадратичного критерію до викидів ($\text{Error} = 59.64$ px).}}
	\end{figure}
	
	\subsection{Частина C. Дослідження власного RANSAC та вибір порога}
	Для власної реалізації RANSAC при робочому порозі $\tau = 3.0$ px отримано:
	\begin{itemize}
		\item Кількість узгоджених точок: 3060 із 3485 ($87.8\%$ інлаєрів);
		\item Спостережувана частка викидів: $\varepsilon = 1 - 3060/3485 \approx 0.122$;
		\item Теоретична кількість ітерацій при $p=0.99$:
		\begin{equation*}
			N = \frac{\ln(1 - 0.99)}{\ln(1 - (1 - 0.122)^4)} = \frac{-4.605}{\ln(1 - 0.596)} = \frac{-4.605}{-0.906} \approx 5.1 \text{ ітерацій}
		\end{equation*}
		\item Середня похибка репроєкції: $0.496$ px.
	\end{itemize}
	
	Звірка з бібліотечною функцією OpenCV (\texttt{cv2.findHomography} з \texttt{cv2.RANSAC}):
	\begin{itemize}
		\item OpenCV RANSAC: 3050 інлаєрів, середня похибка $0.537$ px;
		\item Власний RANSAC: 3060 інлаєрів, середня похибка $0.496$ px.
	\end{itemize}
	Результати практично збігаються, що доводить коректність розробленого алгоритму.
	
	У ході експерименту з варіюванням порога $\tau \in [1, 20]$ px отримано числові дані (табл.~\ref{tab:tau}).
	
	\begin{table}[ht]
		\centering
		\caption{Дослідження впливу порога репроєкції $\tau$}
		\label{tab:tau}
		\begin{tabular}{ccc}
			\toprule
			\textbf{Порог $\tau$ (px)} & \textbf{Кількість інлаєрів} & \textbf{Середня похибка репроєкції (px)} \\
			\midrule
			1 & 2766 & 0.401 \\
			2 & 3004 & 0.477 \\
			3 & 3050 & 0.537 \\
			5 & 3083 & 0.526 \\
			8 & 3092 & 0.532 \\
			12 & 3101 & 0.558 \\
			20 & 3195 & 4.161 \\
			\bottomrule
		\end{tabular}
	\end{table}
	
	\textbf{Обґрунтування робочого порога:} У діапазоні $\tau \in [2, 12]$ px спостерігається плато стабільності (похибка залишається в межах $\approx 0.5$ px). При $\tau = 20$ px похибка стрімко стрибає до $4.161$ px ($+645\%$), що свідчить про хибне включення викидів до множини узгоджених точок. Оптимальним вибором є $\tau = 3.0$ px.
	
	\begin{figure}[ht]
		\centering
		\framebox{\parbox{0.92\textwidth}{\centering \textbf{Рис. 3.} Залежність кількості узгоджених точок та середньої похибки репроєкції від $\tau$ (\texttt{results/step11\_tau\_analysis.png}).\\
				\small Чітке плато в діапазоні $2\text{--}12$ px та різке наростання похибки при $\tau=20$ px через захоплення аутлаєрів.}}
	\end{figure}
	
	\begin{figure}[ht]
		\centering
		\framebox{\parbox{0.92\textwidth}{\centering \textbf{Рис. 4.} Візуалізація інлаєрів та аутлаєрів RANSAC (\texttt{results/step12\_inliers\_outliers.jpg}).\\
				\small Зеленим кольором позначено узгоджені пари на набережній та будівлях; червоним кольором --- відсіяні викиди на воді та хмарах.}}
	\end{figure}
	
	\subsection{Частина D. Зшивання панорами з 5 кадрів}
	Для серії $I_0, I_1, I_2, I_3, I_4$ опорним обрано центральний кадр $I_2$. Перетворення обчислено шляхом композиції матриць гомографій суміжних пар:
	\begin{equation}
		H_{2 \leftarrow 2} = I, \quad H_{1 \leftarrow 2} = H_{1 \to 2}, \quad H_{0 \leftarrow 2} = H_{1 \to 2} \cdot H_{0 \to 1}
	\end{equation}
	\begin{equation}
		H_{3 \leftarrow 2} = (H_{2 \to 3})^{-1}, \quad H_{4 \leftarrow 2} = (H_{2 \to 3})^{-1} \cdot (H_{3 \to 4})^{-1}
	\end{equation}
	
	Знайдено габарити загального полотна через перенесення 4 кутів кожного кадру та враховано зсув $T$ для компенсації від'ємних координат. Накладання виконано із застосуванням двовимірного вагового згладжування (\texttt{distanceTransform}) у зонах перекриття.
	
	Середні похибки репроєкції по всіх парах склали:
	\begin{itemize}
		\item Пара 1--2: $0.537$ px;
		\item Пара 2--3: $0.544$ px;
		\item Пара 3--4: $0.911$ px;
		\item Пара 4--5: $0.919$ px.
	\end{itemize}
	Загальна середня похибка зшивання панорами: $\mathbf{0.728}$ \textbf{px}, що свідчить про високу геометричну узгодженість (субпіксельна точність).
	
	\begin{figure}[ht]
		\centering
		\framebox{\parbox{0.92\textwidth}{\centering \textbf{Рис. 5.} Результат зшивання панорами з 5 послідовних кадрів (\texttt{results/step15\_panorama.jpg}).\\
				\small Повна панорама без швів та артефактів із застосуванням зваженого лінійного згладжування на базі перетворення відстаней.}}
	\end{figure}
	
	\subsection{Частина E. Аналіз випадку відмови}
	Проведено експеримент зіставлення крайніх кадрів серії (Кадр 1 та Кадр 5), перекриття між якими становить $0\%$:
	\begin{itemize}
		\item Знайдено первинних кандидатів: 117 пар;
		\item Після запуску RANSAC узгоджено лише 6 точок;
		\item Фізична модель гомографії відмовила, оскільки вхідні кадри не містять спільної площини спостереження.
	\end{itemize}
	
	% ==========================================
	% ВИСНОВКИ
	% ==========================================
	\section{Висновки}
	\begin{enumerate}
		\item Досліджено роботу SIFT-дескрипторів та показано, що критерій відношення відстаней Лоу ($T = 0.75$) разом із взаємною перевіркою відсіює понад $60\%$ шуму та амбівалентних дескрипторів.
		\item Експериментально доведено непридатність методу найменших квадратів для обчислення гомографії у сирих даних: похибка репроєкції склала $59.64$ px проти $0.496$ px у робастного методу RANSAC.
		\item Успішно розроблено та протестовано власний алгоритм RANSAC. Отримано практичну збіжність із вбудованою реалізацією OpenCV за кількістю інлаєрів ($3060$ проти $3050$).
		\item На основі аналізу залежності метрик від $\tau$ виявлено плато стабільності в межах $\tau \in [2, 12]$ px та обґрунтовано вибір робочого порога $\tau = 3.0$ px.
		\item Зшито цілісну панораму з 5 кадрів через кумулятивне множення гомографій відносно центрального кадру із сумарною середньою похибкою репроєкції $0.728$ px.
	\end{enumerate}
	
\end{document}